\section{Mindset and Motivation in AI Behavior}
\label{sec:1.3}
Three ideas borrowed from how humans develop recur across the design chapters: the stance a system takes toward its own failures, what drives it to act absent an external reward, and what it takes to build any of this with people from diverse disciplines.

\runin{Growth Mindset}
How a system treats its own failures is something the build decides. Carol Dweck's contrast between ability treated as fixed and ability treated as built started as a claim about people; what it asks of a design is whether a system can hold its own mistakes as evidence about itself.

\runin{Intrinsic Motivation and Autonomy}
A moral sense that only activates under supervision is not much of one. Acting well for reasons other than an externally specified reward is the engineering analog of intrinsic motivation, and a system with something like that drive gains independence from moment-to-moment human correction — which leaves open what keeps it aligned absent that correction.

\runin{Interdisciplinary Collaboration}
A machine-learning lab, a psychology department, and an ethics program rarely share a budget, a dataset, or a deadline; building AI this way means changing that default, not hoping the fields converge on their own. What it requires is funding structured to force it and venues that make it a byproduct of doing the work, and then the international partnerships that let it hold across borders.
